% TOP_PROBLEM: {"id":30007633,"problem_number":"LOCAL-30007633","title":"Minimum-wage thresholds and adjustment","category_id":21,"category":{"id":21,"name":"economics","display_name":"Economics","description":"Open economic theory statements, published research agendas, and proposed scoped specifications; question provenance and historical priority are qualified per record.","slug":"economics","order_index":21,"created_at":"2026-10-08T00:00:00Z"},"status":"research_question","external_url":null,"legacy_ids":[],"published":false,"statement":"In the explicitly specified setting: At what minimum-wage levels do gains from correcting employer power give way to employment losses, automation, price increases, or lower job quality? The mathematical statement above fixes the target, assumptions and scope of this catalogue formulation.","economics":{"original_id":122,"field":"Labor markets and work","kind":"empirical","formalization_basis":"adapted_research_question","source_status":"explicit_open","priority_status":"earliest_found","priority_note":"Earliest directly inspected qualifying passage in this audit; earlier origins were not established. A review or research-agenda statement does not imply original priority.","earliest_verified_reference":{"authors":["Patrick Kline"],"title":"Labor Market Monopsony: Fundamentals and Frontiers","year":2025,"url":"https://eml.berkeley.edu/~pkline/papers/HoLE_vol6_monopsony.pdf","doi":"10.1016/bs.heslab.2025.07.007","locator":"Minimum-wage discussion, printed p. 714, paragraph beginning with heterogeneous product demand","evidence_summary":"The quantitative compatibility of joint employment and price responses with monopsony is explicitly posed as future work."},"current_reference":{"authors":["Patrick Kline"],"title":"Labor Market Monopsony: Fundamentals and Frontiers","year":2025,"url":"https://eml.berkeley.edu/~pkline/papers/HoLE_vol6_monopsony.pdf","doi":"10.1016/bs.heslab.2025.07.007","locator":"Minimum-wage discussion, printed p. 714, paragraph beginning with heterogeneous product demand","evidence_summary":"The quantitative compatibility of joint employment and price responses with monopsony is explicitly posed as future work."},"formalization_note":"The cited passage establishes a research direction, not this exact mathematical statement. The notation, population, policy menu, assumptions, and estimands are an original adaptation for this catalogue; no universal unsolved theorem or source priority is claimed."},"statusQualification":"Published question or research agenda verified at the cited locator. The mathematical specification is adapted for this catalogue and includes additional assumptions; source publication does not establish first-ever priority or certify this exact model remains unsolved. The cited passage establishes a research direction, not this exact mathematical statement. The notation, population, policy menu, assumptions, and estimands are an original adaptation for this catalogue; no universal unsolved theorem or source priority is claimed.","statusReviewedAt":"2026-10-08"}
% FIRST_POSTED: 2026-10-08
% ENTRY_KIND: statement_only
% !TeX program = lualatex
\documentclass[11pt]{article}
\usepackage{fontspec}
\usepackage[a4paper,margin=25mm]{geometry}
\usepackage{amsmath,amssymb,mathtools}
\usepackage{xurl}
\usepackage[hidelinks]{hyperref}
\newcommand{\sref}[2]{\hyperlink{source:#1:#2}{[S#2]}}
\setlength{\emergencystretch}{3em}
\begin{document}
% problemId:30007633
\section*{Minimum-wage thresholds and adjustment}\label{30007633}
\subsection{Definitions and mathematical statement}
Fix a set of labor markets, a baseline wage floor \(m_0\), and a policy horizon \(T\). For each feasible minimum wage \(m\), let \(L_t(m)\), \(P_t(m)\), \(A_t(m)\), and \(Q_t(m)\) denote equilibrium employment, output prices, automation investment, and a predeclared job-quality index at date \(t\). These outcomes include employer entry, exit, and worker reallocation. Define \(\Delta_L(m,T)=E[L_T(m)-L_T(m_0)]\), with analogous changes for the other outcomes. Let \(\mathcal M\) be a specified class of models containing employer wage-setting power, heterogeneous product demand, and labor adjustment costs. The research task is to identify joint response curves over the observed policy support and determine whether one model in \(\mathcal M\) can explain their magnitudes simultaneously. If justified by the data, estimate \(m^*=\inf\{m>m_0:\Delta_L(m,T)<0\}\), defining it as infinity if the set is empty. Do not assume a single threshold exists or extrapolate it across countries. Identification must state the policy variation, counterfactual trends, and treatment of spillovers. Distributional welfare conclusions additionally require explicit worker preferences, consumer incidence, and fiscal assumptions.

\par\textbf{Formulation and scope.} The cited passage establishes a research direction, not this exact mathematical statement. The notation, population, policy menu, assumptions, and estimands are an original adaptation for this catalogue; no universal unsolved theorem or source priority is claimed.

\par\textbf{Open-question provenance.} An explicit question or research agenda was verified at the source locator. This does not by itself certify that every later version remains unresolved \sref{30007633}{1}.

\par\textbf{Historical priority.} Earliest directly inspected qualifying passage in this audit; earlier origins were not established. A review or research-agenda statement does not imply original priority.

\subsection{Short English statement}
In the explicitly specified setting: At what minimum-wage levels do gains from correcting employer power give way to employment losses, automation, price increases, or lower job quality? The mathematical statement above fixes the target, assumptions and scope of this catalogue formulation.

\subsection{Sources}
\begin{itemize}\raggedright
\item[\textnormal{[S1]}]\hypertarget{source:30007633:1}{} Patrick Kline. 2025. Labor Market Monopsony: Fundamentals and Frontiers. Minimum-wage discussion, printed p. 714, paragraph beginning with heterogeneous product demand.
\url{https://eml.berkeley.edu/~pkline/papers/HoLE_vol6_monopsony.pdf}.
DOI: 10.1016/bs.heslab.2025.07.007.
{\small\textit{Earliest verified question/agenda reference; historical priority qualified below. The quantitative compatibility of joint employment and price responses with monopsony is explicitly posed as future work.}}
\end{itemize}
\end{document}
